\begin{summarybox}
	\textbf{\textcolor{CSummary}{一句话核心}}
	\textbf{等差数列前$n$项和的最值出现在通项符号由正变负（或由负变正）的时刻，转化于二次函数图象的顶点附近整数点。}
	\medskip
	\textbf{\textcolor{CSummary}{使用条件}}
	\begin{itemize}[leftmargin=*, itemsep=0.5em, topsep=0.4em]
		\item \textbf{条件 1（公差非零）：} 等差数列公差$d\neq0$。
		\item \textbf{条件 2（符号拐点）：} 存在$n$满足$a_n\ge0$且$a_{n+1}\le0$（$d<0$），或$a_n\le0$且$a_{n+1}\ge0$（$d>0$）。
	\end{itemize}
	\medskip
	\textbf{\textcolor{CSummary}{关键提醒}}
	\begin{itemize}[leftmargin=*, itemsep=0.5em, topsep=0.4em]
		\item \textbf{易错点（单侧忽略）：} 只检查$a_n$的符号而忽视$a_{n+1}$，会导致最值点误判。
		\item \textbf{检查项（双值检验）：} 列不等式组时应同时包含$a_n$和$a_{n+1}$，并注意等号可取情形。
	\end{itemize}
\end{summarybox}
